% ==================================================================
% Appendix E — Phase signal-state strings   (writer W5)
% ==================================================================
\chapter{Phase Signal-State Strings}\label{app:phases}

\sumo{} describes the state of a traffic light as a string with one character per controlled
\emph{link}, where a link is a connection from an incoming lane to an outgoing lane through the
junction. The characters used here are: \texttt{G} for a protected green (priority movement);
\texttt{g} for a permissive green that must yield to conflicting traffic; \texttt{y} for amber; and
\texttt{r} for red. \prism{} and all baselines use only protected greens, so \texttt{g} never
appears. Right turns are served only in their own protected phases, and opposing through movements
do not conflict.

The four channelised left-turn slip lanes are \emph{not} part of the traffic-light program. They join
the exit arms through unsignalised yield links and therefore have no character in the strings. Each
arm thus contributes three signalised links, and the junction has twelve. The link order assumed in
this thesis is given in \cref{tab:app-links}, with arms numbered clockwise from the first bypass
arm. \textbf{The actual index order generated by \texttt{netconvert} must be checked in
\texttt{netedit} (traffic-light mode, ``show link indices'') and the strings below re-ordered if
necessary.}

\begin{table}[htbp]
\centering
\caption{Assumed link-index order of the traffic-light program}\label{tab:app-links}
\small
\begin{tabular}{@{}clll@{}}
\toprule
Index & Arm & From lane & Movement\\
\midrule
0 & Bypass 1 & kerb lane & through\\
1 & Bypass 1 & offside lane & through\\
2 & Bypass 1 & offside lane & right\\
3 & Cross 1 & kerb lane & through\\
4 & Cross 1 & offside lane & through\\
5 & Cross 1 & offside lane & right\\
6 & Bypass 2 & kerb lane & through\\
7 & Bypass 2 & offside lane & through\\
8 & Bypass 2 & offside lane & right\\
9 & Cross 2 & kerb lane & through\\
10 & Cross 2 & offside lane & through\\
11 & Cross 2 & offside lane & right\\
\midrule
-- & every arm & slip lane & left (unsignalised yield, not in TLS)\\
\bottomrule
\end{tabular}
\end{table}

\begin{table}[htbp]
\centering
\caption{Signal-state strings for the four phases, intergreens and fail-safe plan}\label{tab:app-states}
\small
\begin{tabular}{@{}llll@{}}
\toprule
Phase & State string (index 0 $\rightarrow$ 11) & Duration & Used by\\
\midrule
$\phi_1$ green (bypass through) & \texttt{GGrrrrGGrrrr} & variable & all controllers\\
$\phi_1$ amber & \texttt{yyrrrryyrrrr} & 3 s & all\\
$\phi_2$ green (bypass right) & \texttt{rrGrrrrrGrrr} & variable & all\\
$\phi_2$ amber & \texttt{rryrrrrryrrr} & 3 s & all\\
$\phi_3$ green (cross through) & \texttt{rrrGGrrrrGGr} & variable & all\\
$\phi_3$ amber & \texttt{rrryyrrrryyr} & 3 s & all\\
$\phi_4$ green (cross right) & \texttt{rrrrrGrrrrrG} & variable & all\\
$\phi_4$ amber & \texttt{rrrrryrrrrry} & 3 s & all\\
All-red (after every amber) & \texttt{rrrrrrrrrrrr} & 2 s & all\\
\bottomrule
\end{tabular}
\end{table}

\Cref{tab:app-failsafe} lists the fail-safe program. It is identical to the observed plan (B1) and is
loaded as the static program \texttt{B1\_fixed78} when the \vls{} delivers no data for five
consecutive scans. Because \prism{} chooses phases acyclically, an amber is always followed by the
all-red state before any other green. The safety test in \cref{app:code} checks this transition rule
and the absence of conflicting \texttt{G} pairs on every simulation step.

\begin{table}[htbp]
\centering
\caption{Fail-safe (B1) program, cycle 78 s}\label{tab:app-failsafe}
\small
\begin{tabular}{@{}cllr@{}}
\toprule
Step & Meaning & State string & Duration (s)\\
\midrule
0 & $\phi_1$ green & \texttt{GGrrrrGGrrrr} & 21\\
1 & amber & \texttt{yyrrrryyrrrr} & 3\\
2 & all-red & \texttt{rrrrrrrrrrrr} & 2\\
3 & $\phi_2$ green & \texttt{rrGrrrrrGrrr} & 8\\
4 & amber & \texttt{rryrrrrryrrr} & 3\\
5 & all-red & \texttt{rrrrrrrrrrrr} & 2\\
6 & $\phi_3$ green & \texttt{rrrGGrrrrGGr} & 15\\
7 & amber & \texttt{rrryyrrrryyr} & 3\\
8 & all-red & \texttt{rrrrrrrrrrrr} & 2\\
9 & $\phi_4$ green & \texttt{rrrrrGrrrrrG} & \est{14}\\
10 & amber & \texttt{rrrrryrrrrry} & 3\\
11 & all-red & \texttt{rrrrrrrrrrrr} & 2\\
\midrule
 & \multicolumn{2}{l}{Cycle} & 78\\
\bottomrule
\end{tabular}
\end{table}
